\documentclass[10pt,a4paper]{amsart}
\usepackage[margin=19mm]{geometry}
\usepackage{amsmath,amssymb,mathtools}
\usepackage[scr=rsfs,cal=euler]{mathalfa}
\usepackage{xfrac}
\usepackage[unicode,hidelinks,bookmarks=false]{hyperref}
\newcommand{\ie}{i.e.\ }
\newcommand{\cf}{cf.\ }
\newcommand{\resp}{respectively\ }
\newcommand{\Subsection}{\subsection}
\providecommand{\nobreakdash}{\nobreak\hbox{-}}
\setlength{\emergencystretch}{2em}
\multlinegap0pt
\usepackage{mathtools}
\usepackage{amssymb}
\usepackage[shortlabels]{enumitem}
\usepackage{xcolor}
\newtheorem{proposition}{Proposition}
\newcommand{\Q}{\mathbb{Q}}
\newcommand{\Z}{\mathbb{Z}}
\newcommand{\mathcalE}{\mathcal{E}}
\title{Torsion and Positive Rank in an Elliptic Family Arising from Cubic 2-Cycles}
\author{Sompong Chuysurichay, Sawian Jaidee,, Chatchawan Panraksa}
\date{Source: arXiv:2606.15109v2 (CC BY 4.0)}
\begin{document}
\maketitle

\noindent\emph{Source and licence.} Torsion and Positive Rank in an Elliptic Family Arising from Cubic 2-Cycles. Version arXiv:2606.15109v2 (CC BY 4.0), distributed by arXiv under the Creative Commons Attribution 4.0 International licence, http://creativecommons.org/licenses/by/4.0/, as recorded in the arXiv licence metadata for this exact version and shown on the abstract page https://arxiv.org/abs/2606.15109v2. Full licence notice: \texttt{licenses/arxiv-2606.15109-CC-BY-4.0.txt}.\par\smallskip

\noindent\textit{Editorial scope.} Counted unit: the article's proposition prop:Et-2torsion (source0.tex lines 551--565). The article's complete original proof of it is delivered with it (source0.tex lines 567--613, one \texttt{proof} environment of 47 lines). No mathematics is written, repaired or re-derived: the statement and the proof are the article's own lines, in the article's own notation, and no retained line is substituted or re-flowed.  Retained uncounted as the source-local context the proof needs: lines 446--448; lines 531--535.  Nothing was omitted from any retained span, so the package carries no omission record.  No retained line contains a citation, so the wrapper carries no bibliography and no bibliography entry.  No retained line contains a figure environment, an image-inclusion command or drawing code, so no image is delivered and the package declares no asset dependency.  Editorial formatting only, in addition: an AMS \texttt{amsart} preamble that restates the article's own declarations and macros used by the retained lines, namely the environments \texttt{proposition} on separate counters, together with the standard packages those lines need.  The retained environments print in this document as Proposition 1 in order of appearance, whereas the article prints the same items with the numbering of its own full document.  The complete native source of the article as distributed in the arXiv e-print for this exact version is bundled unchanged as \texttt{sources/202832/source0.tex}.
\begin{equation}\label{eq:Et-family}
\mathcalE_t:\quad Y^2=X^3+4X^2+16t^2.
\end{equation}

\begin{equation}\label{eq:Et-discriminant}
\Delta(\mathcalE_t)
=-b_2^2b_8-8b_4^3-27b_6^2+9b_2b_4b_6
=-4096t^2(16+27t^2),
\end{equation}

\begin{proposition}\label{prop:Et-2torsion}
Let $t\in\Q^\times$. Then $\mathcalE_t(\Q)$ has a nontrivial rational point of order $2$ if and only if there exists $u\in\Q^\times$ such that
\begin{equation}\label{eq:Et-2torsion-param}
t=\frac{u(u^2+4)}{4}.
\end{equation}
In that case the unique nontrivial rational point of order $2$ is
\begin{equation}\label{eq:Et-2torsion-point}
T_u=\bigl(-(u^2+4),0\bigr),
\end{equation}
and therefore
\[
\mathcalE_t(\Q)[2]\cong \Z/2\Z.
\]
If no such $u$ exists, then $\mathcalE_t(\Q)[2]$ is trivial.
\end{proposition}

\begin{proof}
A rational point on \eqref{eq:Et-family} has order $2$ if and only if it has the form $(r,0)$ with
\begin{equation}\label{eq:Et-2torsion-cubic}
r^3+4r^2+16t^2=0.
\end{equation}
Equivalently,
\begin{equation}\label{eq:Et-2torsion-square}
16t^2=-r^2(r+4).
\end{equation}
Since $t\neq 0$, we also have $r\neq 0$, and \eqref{eq:Et-2torsion-square} gives $-(r+4)=(4t/r)^2$. Write
\[
-(r+4)=u^2,\qquad u\in\Q^\times.
\]
Then
\[
r=-(u^2+4),
\]
and substituting into \eqref{eq:Et-2torsion-square} gives
\[
16t^2=u^2(u^2+4)^2.
\]
Hence
\[
t=\pm \frac{u(u^2+4)}{4}.
\]
Replacing $u$ by $-u$ if necessary yields \eqref{eq:Et-2torsion-param}.

Conversely, if \eqref{eq:Et-2torsion-param} holds, then substituting $X=-(u^2+4)$ into the right-hand side of \eqref{eq:Et-family} gives
\[
X^3+4X^2+16t^2=-(u^2+4)^3+4(u^2+4)^2+u^2(u^2+4)^2=0,
\]
so \eqref{eq:Et-2torsion-point} is indeed a rational point of order $2$.

It remains to prove uniqueness. Writing $\Delta(\mathcalE_t)$ for the
elliptic-curve discriminant in \eqref{eq:Et-discriminant}, the
cubic-polynomial discriminant of $F_t(X):=X^3+4X^2+16t^2$ is
\[
\operatorname{disc}(F_t)=-256t^2(16+27t^2)
=\frac{1}{16}\Delta(\mathcalE_t).
\]
For every $t\in\Q^\times$, this polynomial discriminant is negative.
Therefore the real cubic \(F_t\) has exactly one real root. Hence it has at
most one rational root, and
$\mathcalE_t(\Q)$ has at most one nontrivial rational $2$-torsion point. Thus
the rational $2$-torsion subgroup is either trivial or isomorphic to
$\Z/2\Z$.
\end{proof}

\end{document}
